\documentclass[10pt]{article}
\usepackage[spanish]{babel}
\usepackage[utf8]{inputenc}
\usepackage{graphicx}
\usepackage{hyperref}
\usepackage[lmargin=2cm, rmargin=2cm, top=2cm, bottom=2 cm]{geometry}
\usepackage{fancyhdr}
\pagestyle{fancy}
\setlength{\headheight}{26pt}

\usepackage[table,xcdraw]{xcolor}
\fancyhead{}
\fancyhead[R]{Facultad de Ingeniería y Ciencias \\ Instituto de Ciencias Básicas}
\fancyhead[L]{\includegraphics[width=3.2cm]{img/udp_logo.png}}
\usepackage{animate}
\usepackage{enumerate}
\usepackage{float}
\usepackage{url}
\usepackage{karnaugh-map}
\usepackage{array}
\usepackage{longtable}
\usepackage{pdflscape}

\fancyfoot{}
\fancyfoot[R]{Página \thepage \hspace{0.02 cm}}
\renewcommand{\headrulewidth}{0.9pt}
\renewcommand{\footrulewidth}{0.5pt}

\usepackage{listings}
\usepackage{xcolor}

%New colors defined below
\definecolor{codegreen}{rgb}{0,0.6,0}
\definecolor{codegray}{rgb}{0.5,0.5,0.5}
\definecolor{codepurple}{rgb}{0.58,0,0.82}
\definecolor{backcolour}{rgb}{0.95,0.95,0.92}

%Code listing style named "mystyle"
\lstdefinestyle{mystyle}{
  backgroundcolor=\color{backcolour},   commentstyle=\color{codegreen},
  keywordstyle=\color{magenta},
  numberstyle=\tiny\color{codegray},
  stringstyle=\color{codepurple},
  basicstyle=\ttfamily\footnotesize,
  breakatwhitespace=false,         
  breaklines=true,                 
  captionpos=b,                    
  keepspaces=true,                 
  numbers=left,                    
  numbersep=5pt,                  
  showspaces=false,                
  showstringspaces=false,
  showtabs=false,                  
  tabsize=2
}

\lstset{style=mystyle}

\begin{document}
\pagenumbering{gobble}
\date{}  
\begin{titlepage}
\begin{center}
    \vspace*{\baselineskip}
    
    {
    \bf\fontsize{19}{0}{\selectfont{UNIVERSIDAD DIEGO PORTALES}}\\[0.5cm]
    \fontsize{11}{0}{FACULTAD DE INGENIERÍA Y CIENCIAS}
    }
    
    \vspace*{0.5\baselineskip}
    {
    \bf\fontsize{11}{0}{\selectfont{ESCUELA DE INFORMÁTICA Y TELECOMUNICACIONES}}\\[0.35cm]
    }
    
    \vspace*{\baselineskip}
    \includegraphics[scale=0.50]{img/udp_logo.png}
    \vspace*{3\baselineskip}
     \hrule height 0.5pt
    \vspace{1mm}
    \hrule height 1.5pt
    \vspace*{1\baselineskip}
   
    {
    \bf\fontsize{15}{0}{\selectfont{Software web de gestión ISO \\[0.3cm] Informe Final -- Taller de titulación}}
    }
    \vspace*{1\baselineskip}
    \hrule height 0.5pt
    \vspace{1mm}
    \hrule height 1.5pt
    
    
    \vspace*{4.5\baselineskip}
    
    {
     \bf\fontsize{14}{0}{\selectfont{Profesor:\\[0.3cm]
}}

\bf\fontsize{14}{0}{\selectfont{Jorge Elliott Sotomayor\\[0.5cm]
}}


   

   \bf\fontsize{14}{0}{\selectfont{Estudiante:\\[0.3cm]}}
   \bf\fontsize{14}{0}{\selectfont{Maximiliano Esteban Abascal Fredes}}
   }
    
    \vfill
    Santiago, Chile \hfill 01 de julio de 2026
    
\end{center}
\end{titlepage}
\pagenumbering{roman}
\vspace*{\baselineskip}
\tableofcontents 
\thispagestyle{empty}
\newpage
\pagenumbering{arabic}


\section{Introducción}

Asesorías ISO RCR SpA es una empresa chilena de consultoría y capacitación enfocada en apoyar a organizaciones en la implementación, auditoría y certificación de Sistemas de Gestión bajo normas internacionales ISO. Dentro de su operación actual, la empresa realiza diagnósticos de brecha, acompaña procesos de auditoría y entrega herramientas de control para que sus clientes puedan avanzar hacia la certificación, especialmente bajo la norma ISO 9001:2015 y extendiéndose a otras normativas ISO.

El proyecto presentado en este reporte corresponde al desarrollo de una plataforma web para ejecutar GAP Análisis y gestionar evidencias asociadas a cláusulas de ISO 9001:2015. El GAP Análisis es una auditoría preliminar que compara el estado actual de una organización con los requisitos de una norma, permitiendo identificar incumplimientos, evidencias faltantes y acciones necesarias para avanzar hacia la certificación.

Actualmente, el proceso de diagnóstico y seguimiento se apoya fuertemente en hojas de cálculo, carpetas compartidas y comunicación manual entre consultores, evaluadores y responsables del Sistema de Gestión de Calidad (SGC). Esta forma de trabajo dificulta la trazabilidad, aumenta el riesgo de pérdida o duplicación de información y limita la segregación de datos entre clientes. Por ello, el proyecto busca sentar las bases de un servicio web tipo SaaS que permita estandarizar el diagnóstico, centralizar evidencias y controlar el avance de No Conformidades (NC) de forma segura.

El presente informe consolida el estado final del proyecto, incorporando la información existente del reporte preliminar de junio, el Documento de Especificación de Requisitos de Software (SRS) versión 3.1, la hoja de seguimiento de tareas, los ajustes de alcance definidos para el MVP y la documentación técnica generada durante el desarrollo: 11 Architecture Decision Records (ADR), registro de deuda técnica, diagramas C4 detallados en Mermaid, documento de componentes críticos y plan de pruebas E2E. El proyecto se ejecutó en una planificación de seis semanas de desarrollo tras el trabajo inicial de levantamiento de requerimientos y diseño, totalizando aproximadamente 210 horas de implementación.

\newpage
\section{Antecedentes de la empresa}

Asesorías ISO RCR SpA presta servicios de consultoría, capacitación y acompañamiento a organizaciones que requieren implementar o mejorar sus Sistemas de Gestión. Su propuesta de valor se relaciona con traducir requisitos normativos en planes de trabajo prácticos, auditables y sostenibles para cada cliente. En este contexto, la digitalización del GAP Análisis no solo mejora la operación interna de la consultora, sino que también habilita una oferta de servicio más escalable y trazable.

El rubro de la consultoría ISO requiere mantener consistencia metodológica, control documental y evidencia verificable. En particular, ISO 9001:2015 exige demostrar cumplimiento mediante registros, documentos, responsabilidades y seguimiento de acciones. Cuando estos elementos se administran manualmente, la consultora depende de archivos dispersos y controles externos que no siempre aseguran integridad, versionamiento ni separación adecuada entre clientes.

La estrategia comercial y metodológica actual de Asesorías ISO RCR SpA se apoya en la metodología \textit{Sprint ISO}, presentada por la empresa como un camino desde la planificación hasta la certificación en pasos simples. Esta metodología considera seis etapas: diagnóstico inicial, planificación, implementación, capacitación, auditoría interna y certificación. El proyecto se alinea con esta forma de trabajo porque digitaliza componentes centrales del servicio: diagnóstico GAP, gestión de evidencias, responsables, No Conformidades y trazabilidad de avance.

\begin{table}[H]
\centering
\begin{tabular}{|p{3.2cm}|p{6.1cm}|p{5.1cm}|}
\hline
\rowcolor{gray!20}\textbf{Etapa Sprint ISO} & \textbf{Propósito en la metodología de la empresa} & \textbf{Relación con el MVP} \\ \hline
Diagnóstico inicial & Evaluar brechas mediante auditoría GAP y definir un alcance acotado. & Motor GAP, requisitos ISO y registro de cumplimiento. \\ \hline
Planificación & Definir tiempos, responsables y documentación requerida. & Espacios de trabajo, usuarios, roles y seguimiento de estados. \\ \hline
Implementación & Aplicar procesos y herramientas prácticas en la operación del cliente. & Gestión de evidencias, acciones correctivas y seguimiento de No Conformidades. \\ \hline
Capacitación & Preparar al equipo del cliente para adoptar la norma y mantener buenas prácticas. & Interfaz centralizada y estructura comprensible para responsables del SGC. \\ \hline
Auditoría interna & Revisar cumplimiento y corregir hallazgos antes de la certificación. & Revisión de evidencias, validación, historial y estados de No Conformidades. \\ \hline
Certificación & Acompañar la auditoría externa y respaldar el cumplimiento. & Trazabilidad documental y consolidación de información para auditoría. \\ \hline
\end{tabular}
\caption{Relación entre la metodología Sprint ISO y el alcance del MVP.}
\label{tab:sprint-iso-mvp}
\end{table}

El proyecto impacta directamente en la operación de Asesorías ISO RCR SpA, ya que transforma un proceso de diagnóstico y seguimiento actualmente manual en un flujo web centralizado. A su vez, afecta positivamente a los clientes, quienes podrán revisar requisitos, evidencias, estados y observaciones dentro de un espacio de trabajo propio.

\subsection{Equipo de trabajo}

El equipo de trabajo del proyecto está compuesto por Maximiliano Esteban Abascal Fredes, quien cumple el rol de diseñador del sistema y desarrollador lead del MVP; Ricardo Candia, CEO de Asesorías ISO RCR SpA, quien actúa como contraparte ejecutiva y orienta la visión de negocio del producto; y María, supervisora del proyecto y consultora especialista en sistemas de gestión ISO, quien apoya la validación funcional, normativa y metodológica de los flujos asociados al GAP Análisis, evidencias y No Conformidades. Esta composición permite combinar desarrollo técnico, dirección estratégica y conocimiento especializado en sistemas de gestión ISO.

\section{Formulación del problema}

El problema central consiste en que la gestión actual de diagnósticos GAP, evidencias y No Conformidades se realiza mediante herramientas fragmentadas que no fueron diseñadas para asegurar trazabilidad normativa, segregación multi-cliente ni control de estados. Esta situación genera cuatro dificultades principales:

\begin{enumerate}
    \item \textbf{Baja trazabilidad:} las decisiones, cambios de estado y reemplazos de evidencias pueden quedar repartidos entre planillas, correos o conversaciones externas.
    \item \textbf{Riesgo de mezcla de información:} al trabajar con múltiples clientes, es necesario asegurar que cada organización acceda únicamente a sus propios datos.
    \item \textbf{Seguimiento manual de No Conformidades (NC):} las No Conformidades requieren pasar por estados definidos, comentarios, revisión y cierre, pero el control manual dificulta conocer el avance real.
    \item \textbf{Falta de vista centralizada de evidencias:} no existe una vista única que permita revisar, comparar y validar rápidamente las evidencias asociadas a requisitos, clientes o No Conformidades.
    \item \textbf{Escalabilidad limitada:} mientras más clientes y auditorías existan, mayor es el costo operativo de mantener archivos, versiones y validaciones sin una plataforma común.
\end{enumerate}

La relevancia del proyecto para la empresa radica en que el sistema propuesto reduce errores operacionales, estandariza la metodología de diagnóstico y genera una base tecnológica comercializable. El alcance se concentra en un MVP para ISO 9001:2015, priorizando el núcleo de valor: motor de evaluación GAP, gestión de requisitos, evidencias, usuarios, espacios de trabajo y flujo básico de No Conformidades.

\section{Objetivos}

\subsection{Objetivo general}

Diseñar e implementar un MVP de plataforma web para Asesorías ISO RCR SpA que permita ejecutar GAP Análisis de ISO 9001:2015, gestionar evidencias por requisito y controlar el flujo de No Conformidades mediante roles, espacios de trabajo y trazabilidad básica.

\subsection{Objetivos específicos}

\begin{enumerate}
    \item Definir la arquitectura funcional del sistema a partir del SRS, separando roles, permisos y espacios de trabajo por cliente.
    \item Implementar autenticación, autorización y administración CRUD de usuarios y espacios de trabajo.
    \item Cargar la estructura de requisitos de ISO 9001:2015 mediante seeding de base de datos para evitar cargas manuales en el MVP.
    \item Desarrollar el motor GAP para registrar cumplimiento, abrir No Conformidades, gestionar Acciones Correctivas y calcular indicadores de avance.
    \item Implementar la gestión de evidencias asociadas a requisitos, incluyendo carga, visualización, revisión e historial de versiones.
    \item Validar el sistema mediante pruebas unitarias, pruebas de integración, pruebas de aislamiento multi-tenancy y pruebas UAT con la contraparte.
\end{enumerate}

\begin{table}[H]
\centering
\begin{tabular}{|p{7.2cm}|p{7.2cm}|}
\hline
\rowcolor{gray!20}\textbf{Objetivo específico} & \textbf{Evidencia o estado de cumplimiento en el informe} \\ \hline
Definir arquitectura funcional, roles, permisos y espacios por cliente. & SRS v3.1, definición de roles, modelo multi-tenancy y diagramas C4. \\ \hline
Implementar autenticación, autorización y CRUD de usuarios y espacios. & Módulos implementados de login, JWT, usuarios, roles y espacios de trabajo. \\ \hline
Cargar requisitos ISO 9001:2015 mediante seeding. & Población inicial de requisitos ISO mediante scripts de seeding. \\ \hline
Desarrollar motor GAP, No Conformidades, Acciones Correctivas e indicadores. & Motor GAP funcional, apertura y seguimiento de No Conformidades, acciones correctivas y cálculo de cumplimiento. \\ \hline
Implementar gestión de evidencias con revisión e historial. & Carga, visualización, revisión, vinculación e historial de evidencias. \\ \hline
Validar mediante pruebas unitarias, integración, multi-tenancy y UAT. & Pruebas unitarias e integración ejecutadas; pruebas multi-tenancy y UAT pendientes según cronograma. \\ \hline
\end{tabular}
\caption{Relación entre objetivos específicos y evidencia de avance.}
\label{tab:objetivos-evidencia}
\end{table}

\section{Marco conceptual}

\subsection{ISO 9001:2015 y Sistema de Gestión de Calidad}

ISO 9001:2015 es una norma internacional orientada a Sistemas de Gestión de Calidad. Su aplicación exige que la organización demuestre capacidad para entregar productos o servicios conformes, gestionar riesgos, controlar procesos y mejorar continuamente. En el contexto del proyecto, la norma funciona como estructura base para evaluar requisitos, registrar evidencias y determinar brechas de cumplimiento.

\subsection{GAP Análisis}

El GAP Análisis es un diagnóstico comparativo entre la situación actual de una organización y el estado requerido por una norma. Su resultado permite clasificar requisitos, identificar brechas, priorizar acciones y generar evidencia para futuras auditorías. En la plataforma, este concepto se traduce en un motor que permite revisar cláusulas, registrar estados de cumplimiento y abrir No Conformidades cuando corresponde junto a sus Acciones Correctivas asociadas según sea el caso.

\subsection{No Conformidad}

Una No Conformidad corresponde a una desviación respecto de un requisito normativo o procedimiento definido. Para el MVP, su flujo se simplifica a estados operativos: Abierta, Análisis, Implementación, Verificación y Cerrada. Esta simplificación reemplaza modelos jerárquicos de acciones correctivas padre/hijo y permite concentrar el desarrollo en un seguimiento claro, verificable y alcanzable. También se cuenta con un estado de aceptación de la No Conformidad por parte del Responsable del SGC, cuyo objetivo es confirmar que la desviación fue recibida, comprendida y asumida por el área responsable antes de iniciar el análisis de causa y la implementación de acciones correctivas.

\textbf{Nota terminológica:} En la interfaz de usuario del sistema se utiliza el término ``Brecha Detectada en GAP Analysis'' o simplemente ``Brecha'' en lugar de ``No Conformidad'', con el propósito de hacer más comprensible el concepto para clientes que no poseen experiencia previa con terminología ISO. El presente informe mantiene el término ``No Conformidad'' por ser el vocablo normativo estándar definido en ISO 9001:2015.

\subsection{Multi-tenancy y segregación de datos}

La plataforma se diseña para operar con múltiples clientes dentro de un mismo sistema. Por ello, cada espacio de trabajo representa a una organización cliente y debe aislar usuarios, requisitos, evidencias, No Conformidades, Acciones Correctivas e historial. La segregación de datos es un requisito crítico, ya que evita accesos cruzados y mantiene confidencialidad entre clientes.

\subsection{MVP y reducción de alcance}

El Producto Mínimo Viable se entiende como la versión inicial que entrega el núcleo de valor con la menor complejidad posible. En este proyecto, el MVP se acota a ISO 9001:2015, gestión GAP, evidencias, NC (y sus Acciones Correctivas) y roles esenciales. La reducción de alcance elimina funcionalidades de alta complejidad técnica o bajo valor inmediato para asegurar una entrega sólida en 210 horas.


\subsection{Soluciones y metodologías existentes}

Actualmente existen alternativas para abordar la gestión de auditorías, cumplimiento y sistemas de gestión ISO. La primera alternativa, utilizada de forma extendida en operaciones pequeñas o en etapas iniciales, corresponde al uso de planillas Excel, carpetas compartidas en Drive y comunicación manual. Su ventaja principal es el bajo costo inicial y la flexibilidad inmediata; sin embargo, presenta debilidades importantes para el problema tratado: baja trazabilidad, riesgo de duplicidad, dificultad para controlar versiones, poca segregación entre clientes y dependencia de disciplina manual para mantener los registros actualizados.

Una segunda alternativa corresponde a plataformas comerciales especializadas en gestión ISO, auditorías o GRC, tales como ISWO, GOAudits y GRCTools. De acuerdo con la entrevista de validación realizada con María, ISWO es la alternativa que más se aproxima a la visión general de Asesorías ISO RCR SpA. No obstante, también se identificó que estas plataformas pueden resultar complejas para empresas que no tienen experiencia previa con sistemas ISO, ya que requieren comprender metodologías, parametrizaciones y flujos de trabajo antes de obtener valor operativo. En el caso de ISWO, la principal limitación observada es que ciertos procesos requieren parametrizaciones previas engorrosas para acciones que deberían ser simples, como trabajar una matriz FODA. Además, su IA se percibe como una asistencia conversacional que responde preguntas, pero no opera como agente capaz de modificar, completar o actuar sobre los datos del cliente dentro de la plataforma.

La tercera alternativa corresponde al desarrollo de una plataforma propia. Aunque el MVP actual es más acotado que las soluciones comerciales existentes, permite construir una base ajustada al flujo específico de Asesorías ISO RCR SpA. El valor estratégico de esta alternativa, validado en entrevista, se relaciona con proteger la información de los clientes, mantener control sobre el sistema, evitar el uso de múltiples herramientas dispersas y entregar un servicio de calidad desde una única plataforma centralizada. Además, permite diseñar una experiencia más simple y entendible para clientes que no dominan la terminología ni la metodología ISO.

El diferenciador del producto completo no se limita a reemplazar planillas o replicar software comercial. La visión final apunta a una plataforma integrada, fácilmente auditable, con historial detallado, funciones inteligentes, IA agente, importación inteligente de datos, reutilización de información entre módulos para evitar registros duplicados, relación entre elementos del sistema y escalabilidad hacia otras normas ISO. En este sentido, el MVP se interpreta como la primera base operativa: es más simple, centralizado y hecho a medida, mientras que el producto final busca convertirse en una plataforma inteligente de gestión ISO adaptable a la metodología de la consultora.

\begin{table}[H]
\centering
\begin{tabular}{|p{3.2cm}|p{4cm}|p{4cm}|p{3.4cm}|}
\hline
\rowcolor{gray!20}\textbf{Alternativa} & \textbf{Ventajas} & \textbf{Limitaciones} & \textbf{Encaje con el proyecto} \\ \hline
Excel/Drive & Bajo costo inicial, rapidez de uso y alta flexibilidad manual. & Baja trazabilidad, duplicidad de registros, control documental débil y riesgo de mezcla entre clientes. & Útil como punto de partida, pero insuficiente para centralizar, auditar y escalar el servicio. \\ \hline
ISWO, GOAudits y GRCTools & Plataformas maduras para auditoría, cumplimiento, reportabilidad y gestión de procesos. & Mayor complejidad metodológica, parametrización previa y menor adaptación al flujo propio de Asesorías ISO RCR SpA. & Referentes del estado del arte, pero menos alineados con una experiencia simple y hecha a medida para clientes sin dominio ISO. \\ \hline
Plataforma propia & Control de datos, centralización, simplicidad, auditabilidad, roadmap propio y adaptación a la consultora. & Mayor riesgo técnico inicial y menor cobertura funcional en el MVP. & Mejor alternativa estratégica para evolucionar hacia IA agente, importación inteligente, reutilización de datos y soporte multi-norma. \\ \hline
\end{tabular}
\caption{Comparación de alternativas para la gestión GAP e ISO.}
\label{tab:alternativas}
\end{table}

\subsection{Producto completo proyectado}

El producto completo se proyecta como una plataforma SaaS de gestión ISO que integre, en un mismo entorno, los procesos principales del Sistema de Gestión: contexto de la organización, planificación, liderazgo, apoyo, operación, evaluación y mejora continua. Según la entrevista, estos ejes deben constituir la estructura funcional de largo plazo, mientras que módulos como productividad, mantenimiento, planes de mejoramiento, PQRSF, proveedores, SGSST o dashboards genéricos no forman parte prioritaria de la propuesta final. La visualización analítica sí se considera relevante, pero enfocada en indicadores útiles para el seguimiento del sistema y no como un conjunto amplio de métricas desconectadas.

El diseño completo debe operar cláusula a cláusula, de manera que requisitos, evidencias, hallazgos, No Conformidades, acciones correctivas, aprobaciones e historial queden relacionados desde el origen. Esta estructura busca evitar duplicidad de información, reducir registros repetidos y permitir que una evidencia o acción tenga impacto en los componentes normativos correspondientes sin tener que ser ingresada varias veces. La trazabilidad esperada incluye quién crea un registro, quién firma, quién aprueba, qué cambios se realizaron, cuál era el estado anterior, cuál es el estado posterior y la fecha de cada modificación.

En términos de escalabilidad normativa, la prioridad futura identificada es avanzar desde ISO 9001:2015 hacia ISO 27001, seguridad en el trabajo y medioambiente, incorporando progresivamente normas adicionales bajo una arquitectura reutilizable. La inteligencia artificial queda como una línea estratégica pendiente de evaluación detallada, pero la visión validada apunta a una IA agente capaz de asistir o ejecutar acciones sobre la plataforma, no solo responder preguntas. Esto permitiría apoyar generación de informes, sugerencia de brechas, clasificación de evidencias e importación inteligente de datos en etapas posteriores.

\section{Diseño metodológico y alcance del MVP}

La metodología seguida combina especificación de requisitos, diseño de arquitectura, desarrollo iterativo y validación incremental. El SRS v3.1 actúa como documento base para definir actores, requisitos funcionales, requisitos no funcionales, criterios de validación y exclusiones del MVP. A partir de este diseño, la hoja de seguimiento organiza las tareas por módulos, responsables, fechas, horas estimadas y estado.

\subsection{Metodología de trabajo y seguimiento}

En una primera etapa, el proyecto se gestionó mediante una metodología ágil acompañada por una Carta Gantt. Esta combinación permitió ordenar el levantamiento inicial, estimar tiempos y visualizar dependencias generales entre actividades. Sin embargo, durante la ejecución se evidenció que la Carta Gantt introducía un nivel de rigidez poco conveniente para el contexto real del proyecto, especialmente por el tiempo requerido para registrar avances, reprogramar tareas y ajustar continuamente el plan frente a cambios de alcance o descubrimientos técnicos.

Debido a lo anterior, desde el 6 de mayo se optó por una metodología completamente ágil, apoyada principalmente en una tabla Kanban. Este cambio permitió organizar el trabajo por flujo de estados, priorizar tareas de forma más flexible y adaptar el desarrollo a los hallazgos del MVP sin perder visibilidad sobre lo pendiente, lo activo y lo finalizado. La tabla Kanban facilitó una gestión más liviana y práctica, alineada con el tiempo limitado disponible para construir el núcleo funcional.

A partir de finales de mayo se comenzó a traspasar progresivamente las tareas relevantes a Jira, con el objetivo de profesionalizar el seguimiento, centralizar incidencias y aprender de manera continua el uso de la herramienta. En paralelo, se mantuvo la hoja de cálculo como registro adicional de respaldo, principalmente para conservar una evidencia simple del progreso y contar con una fuente alternativa en caso de requerir comparación o reconstrucción del avance. La proyección es que, en etapas posteriores, el cálculo de progreso y la gestión completa de tareas se realicen directamente en Jira, dejando la planilla solo como apoyo documental si fuese necesario.

\subsection{Justificación de herramientas técnicas}

La selección tecnológica del MVP busca equilibrar rapidez de desarrollo, mantenibilidad, escalabilidad y disponibilidad de documentación oficial. En la capa de presentación se utiliza una SPA construida con Vite y React, coherente con la estructura del repositorio en la carpeta \texttt{frontend/}. React aporta un modelo de componentes reutilizables que facilita dividir la interfaz en piezas mantenibles y reducir deuda técnica, mientras que Vite entrega un entorno de desarrollo liviano, rápido y adecuado para construir una aplicación web de una sola página. Se incorporan además \texttt{chart.js} con \texttt{react-chartjs-2} para gráficos radar de cumplimiento por cláusula en los dashboards, \texttt{react-hook-form} para formularios sin re-renders innecesarios y \texttt{react-idle-timer} para cierre automático de sesión tras inactividad prolongada. En la capa backend se utiliza Express.js sobre Node.js, organizado en la carpeta \texttt{backend-js/}, lo que permite exponer servicios HTTP de forma simple y mantener una separación clara entre frontend, API y persistencia.

La capa de servicios se organiza mediante una API REST, debido a que este estilo arquitectónico desacopla cliente y servidor usando convenciones HTTP y JSON. Esto permite que el frontend, futuras integraciones, automatizaciones o servicios de IA puedan interactuar con el backend sin depender de una interfaz única. Para persistencia se utiliza MySQL, ya que la empresa cuenta con su propio servidor MySQL y el dominio del sistema contiene relaciones claras entre usuarios, espacios de trabajo, requisitos, evidencias, No Conformidades, acciones correctivas e historial. El enfoque relacional permite mantener integridad y consistencia transaccional mediante propiedades ACID, especialmente relevantes para trazabilidad y auditoría (se considera a futuro aplicar una tecnologia de base de datos distinta y mas robusta). Para el almacenamiento de archivos de evidencia se integra Google Drive mediante su API REST y la librería \texttt{googleapis}, organizando carpetas por workspace y ambiente (Development/Production); cada evidencia en la base de datos almacena una referencia \texttt{drive://\{id\}} que permite recuperar el archivo desde Drive sin duplicar contenido en la base relacional.

Para las interacciones en tiempo real se utiliza WebSockets mediante la librería \texttt{ws} integrada directamente en el servidor HTTP de Node.js. Esta decisión reconsidera el ajuste de alcance inicial, donde el chat había quedado oficialmente excluido como WebSocket o Socket.io y se planteaba como JSON persistido y actualizado vía CRUD. Durante el desarrollo se identificó que ese enfoque limitaba la inmediatez de las respuestas y la escalabilidad futura de las conversaciones, ya que los mensajes deben poder asociarse a distintos elementos del dominio, tales como evidencias, No Conformidades, acciones correctivas y futuros módulos. Para las notificaciones internas se reutiliza la misma conexión WebSocket del chat, simplificando la arquitectura y reduciendo la sobrecarga de conexiones simultáneas. El sistema cuenta además con un worker de notificaciones programadas que cada 30 segundos convierte registros de \texttt{SCHEDULED\_NOTIFICATIONS} en notificaciones reales cuando llega su fecha de activación.

En seguridad se adopta JWT para autenticación basada en tokens firmados, siguiendo el estándar RFC 7519. Esta decisión facilita un esquema de autenticación stateless, adecuado para servicios web escalables y separación entre frontend y backend. Para reproducibilidad del entorno se incorporan scripts de seeding, con el fin de cargar la estructura ISO 9001:2015 de forma consistente en desarrollo, pruebas y despliegue. Finalmente, la gestión del proyecto se apoya en Kanban y Jira para limitar trabajo en progreso, visualizar cuellos de botella, mantener trazabilidad de tareas y migrar progresivamente desde la hoja de cálculo hacia una herramienta profesional de gestión ágil.

\subsection{Arquitectura C4}

Para complementar el diseño técnico del MVP se elaboraron diagramas bajo el modelo C4, con el propósito de representar el sistema desde distintos niveles de abstracción. Estos documentos permiten explicar la relación entre usuarios, plataforma, contenedores principales y componentes internos, facilitando la comunicación del diseño entre el equipo técnico y la contraparte de negocio.

El diagrama de contexto presenta la plataforma como un sistema web utilizado por el Administrador, el Evaluador y el Responsable del SGC, destacando su interacción con los espacios de trabajo y la base de datos del servicio.

\begin{figure}[H]
    \centering
    \includegraphics[width=0.92\textwidth]{img/C4-Contexto-fixed-1-2026-06-04-234518.png}
    \caption{Diagrama C4 de contexto del sistema.}
    \label{fig:c4-contexto}
\end{figure}

El diagrama de contenedores resume la separación entre frontend, backend API y base de datos MySQL. Además, el backend considera un canal WebSocket para chat contextual y notificaciones internas en tiempo real, reutilizando la misma conexión para reducir complejidad operativa. Esta vista permite visualizar cómo se distribuyen las responsabilidades principales del MVP y cómo se mantiene la segregación de datos por espacios de trabajo.

\begin{figure}[H]
    \centering
    \includegraphics[width=0.92\textwidth]{img/C4-Contenedor-fix-1-Página-2.png}
    \caption{Diagrama C4 de contenedores del sistema.}
    \label{fig:c4-contenedores}
\end{figure}

A nivel de componentes, se documentaron por separado el frontend y la API backend. El diagrama de frontend muestra el sistema de componentes visuales, enrutamiento y renderizacion de paginas, flujos de navegación y módulos de interacción del usuario, mientras que el diagrama de backend describe los componentes responsables de autenticación, autorización, gestión GAP, evidencias, No Conformidades, historial y persistencia.

\begin{figure}[H]
    \centering
    \includegraphics[width=0.92\textwidth]{img/C4-Componentes-Frontend-2026-06-05-030116.png}
    \caption{Diagrama C4 de componentes del frontend.}
    \label{fig:c4-frontend}
\end{figure}

\begin{figure}[H]
    \centering
    \includegraphics[width=0.78\textwidth]{img/C4-Componentes-BackendAPI-2026-06-05-030115.png}
    \caption{Diagrama C4 de componentes de la API backend.}
    \label{fig:c4-backend}
\end{figure}

\subsection{Roles del sistema}

\begin{table}[H]
\centering
\begin{tabular}{|p{3.2cm}|p{11.2cm}|}
\hline
\rowcolor{gray!20}\textbf{Rol} & \textbf{Responsabilidad principal} \\ \hline
Administrador & Gestiona espacios de trabajo, usuarios, roles y configuración inicial del sistema. Asigna contraseñas genéricas y controla la gobernanza general. \\ \hline
Evaluador & Revisa auditorías, evidencias y No Conformidades (brechas). Puede aprobar, rechazar, comentar y mover estados de revisión sin editar el contenido original. Marca requisitos como No Aplica según ISO 9001 Req 4.3. \\ \hline
Responsable del SGC & Opera el espacio de su organización, registra información, carga evidencias, responde NC, gestiona acciones correctivas y mantiene actualizados los antecedentes de cumplimiento. \\ \hline
Operativo & Visualiza el dashboard de cumplimiento de su workspace, consulta requisitos ISO y revisa el estado de No Conformidades y acciones correctivas. Es un rol de solo lectura orientado a mantener visibilidad operativa sin capacidad de modificación. \\ \hline
\end{tabular}
\caption{Roles principales definidos para el MVP.}
\label{tab:roles}
\end{table}

\subsection{Funcionalidades incluidas}

El alcance vigente del MVP incluye los siguientes módulos:

\begin{enumerate}
    \item Administración de usuarios, roles y espacios de trabajo.
    \item Login seguro con JWT, refresh token y cambio de contraseña dentro de la plataforma.
    \item Activación de cuenta en primer ingreso con cambio obligatorio de contraseña.
    \item Cierre automático de sesión por inactividad prolongada (\texttt{react-idle-timer}).
    \item Seeding de requisitos de ISO 9001:2015 en la base de datos mediante script idempotente.
    \item Navegación por requisitos ISO con árbol jerárquico y apertura de No Conformidades (brechas).
    \item Gestión de estados de NC: Abierta, Análisis, Ejecución, Verificación y Cerrada.
    \item Marcado de requisitos como No Aplica (NA) según ISO 9001:2015 Req 4.3, con exclusión automática del cálculo de cumplimiento.
    \item Gestión de acciones correctivas con estructura de lista enlazada (padre/hijo), tablero Kanban con drag-and-drop, estados (Pendiente, En Progreso, Eficaz, No Eficaz), historial campo por campo y eliminación en cascada.
    \item Carga, visualización, aprobación, rechazo e historial de evidencias almacenadas en Google Drive con carpetas segregadas por workspace.
    \item Chat contextual en tiempo real mediante WebSockets, asociado a No Conformidades, evidencias, acciones correctivas y requisitos.
    \item Notificaciones internas en tiempo real mediante WebSockets, reutilizando la misma conexión del chat, con worker de notificaciones programadas para verificación de eficacia.
    \item Dashboard analítico con gráficos radar de cumplimiento por cláusula, KPIs de resolución y eficiencia, y vistas diferenciadas por rol (Administrador, Evaluador, Responsable SGC, Operativo).
    \item Lobbies diferenciados por rol con métricas y accesos contextualizados.
    \item Pruebas unitarias, de integración y de partición de equivalencia; pruebas multi-tenancy y UAT planificadas.
\end{enumerate}

\subsection{Funcionalidades excluidas por ajuste de alcance}

Para viabilizar el proyecto en seis semanas, se excluyen explícitamente de la Etapa 1 las siguientes capacidades:

\begin{enumerate}
    \item Recuperación de contraseña, onboarding y notificaciones mediante correo electrónico.
    \item WebSockets o Socket.io para chat en tiempo real; el chat queda como JSON persistido y actualizado vía CRUD (exclusión oficial del ajuste de alcance inicial; posteriormente se reconsidera e implementa mediante WebSockets, ya que se requiere inmediatez y asociación escalable de mensajes con evidencias, No Conformidades, acciones correctivas y futuros elementos del sistema).
    \item Almacenamiento en AWS S3; las evidencias se almacenan en Google Drive como solución intermedia más simple, con carpetas segregadas por workspace y referencias \texttt{drive://\{id\}} en la base de datos.
    \item Carga dinámica de requisitos normativos por CSV; ISO 9001:2015 se carga por script de seeding.
    \item Edición documental en línea e importación flexible de herramientas Excel.
    \item Flujos de aprobación complejos, firma múltiple o evaluadores consecutivos.
    \item Árbol jerárquico de acciones correctivas padre/hijo (exclusión reconsiderada: se implementó una lista enlazada donde cada acción referencia a su predecesora mediante \texttt{accion\_previa\_id}, permitiendo profundidad ilimitada y visualización en tablero Kanban).
    \item Historial microscópico de cambios de texto en todas las entidades del sistema; se implementó historial detallado para acciones correctivas (campo por campo) y para estados de No Conformidades y evaluaciones, pero no para todos los campos de texto de la plataforma.
    \item Soporte multi-norma, automatizaciones de riesgo/FODA, PEST, Cinco Fuerzas de Porter y cualquier herramienta no incluida en el flujo GAP/NC básico.
\end{enumerate}

\section{Estado actual de avance}

De acuerdo con la hoja de seguimiento actualizada a la fecha de cierre, la planificación general contempla 281 unidades de trabajo registradas para el MVP vigente. Respecto del reporte preliminar, se han completado módulos significativos que antes figuraban como pendientes, incluyendo dashboard analítico con gráficos radar, tablero Kanban de acciones correctivas, integración con Google Drive para evidencias, worker de notificaciones programadas, lobbies diferenciados por rol y timeout por inactividad.

\subsection{Resumen cuantitativo}

\begin{table}[H]
\centering
\begin{tabular}{|p{4cm}|p{3cm}|p{3cm}|p{5cm}|}
\hline
\rowcolor{gray!20}\textbf{Estado} & \textbf{Total} & \textbf{Porcentaje} & \textbf{Observación} \\ \hline
Completado & 245 & 87,2\% & Incluye todos los módulos del núcleo funcional: autenticación, GAP, NC, acciones correctivas, evidencias, dashboard, Kanban, WebSockets y validaciones base. \\ \hline
En progreso & 12 & 4,3\% & Incluye ajustes de seguridad IDOR, refactoring de componentes y documentación técnica final. \\ \hline
Pendiente & 24 & 8,5\% & Incluye pruebas multi-tenancy formales, pruebas UAT, correcciones de deuda técnica menor y cierre documental. \\ \hline
Total & 281 & 100\% & Plan de trabajo completo registrado en la hoja de avance. \\ \hline
\end{tabular}
\caption{Estado actualizado de unidades de trabajo según hoja de seguimiento del proyecto.}
\label{tab:estado}
\end{table}

En términos de seguimiento actualizado, la hoja registra 245 unidades completadas, 12 en progreso y 24 pendientes. Esto representa un 87,2\% de avance completado sobre las 281 unidades totales, o un 91,5\% si se consideran como avance las unidades completadas más las actualmente en progreso. Para efectos de gestión, el indicador más conservador es el 87,2\%, ya que solo contabiliza trabajo finalizado. Este incremento de aproximadamente 18 puntos porcentuales respecto al 69,4\% reportado en junio refleja la incorporación de los módulos de dashboard, Kanban, integración Google Drive, evaluaciones NA y notificaciones programadas, todos los cuales se encontraban pendientes en el informe anterior.

Como complemento visual, se incorporan tres documentos de seguimiento: el cronograma general de tareas, el gráfico de distribución de estados y la hoja de avance exportada desde la planilla de control. Estos elementos permiten contrastar el avance textual con evidencia gráfica y documental del seguimiento realizado.

\begin{figure}[H]
    \centering
    \includegraphics[width=0.95\textwidth]{img/cronograma_de_tareas_proyecto_iso.png}
    \caption{Cronograma de tareas del Proyecto ISO.}
    \label{fig:cronograma-tareas}
\end{figure}

\begin{figure}[H]
    \centering
    \includegraphics[width=0.75\textwidth]{img/grafico_pastel_proyecto_iso.png}
    \caption{Distribución de tareas por estado del Proyecto ISO.}
    \label{fig:grafico-pastel-avance}
\end{figure}

\begin{figure}[H]
    \centering
    \includegraphics[width=0.95\textwidth]{img/Avance Proyecto ISO - Hoja 1.pdf}
    \caption{Hoja de seguimiento de tareas del Proyecto ISO.}
    \label{fig:hoja-avance-proyecto}
\end{figure}

\subsection{Módulos completados}

Los avances consolidados son los siguientes:

\begin{enumerate}
    \item Implementación de base de datos y arquitectura base.
    \item CRUD de espacios de trabajo con protección de rutas.
    \item Gestión de usuarios, roles, contraseñas y activación de cuenta en primer ingreso.
    \item Autenticación, autorización, persistencia de sesión y redirección.
    \item Timeout por inactividad con cierre automático de sesión.
    \item Sistema de notificaciones internas en tiempo real mediante WebSockets.
    \item Worker de notificaciones programadas para verificación de eficacia.
    \item Población de datos de ISO 9001:2015 mediante seeding.
    \item Interfaz de navegación e interfaz de requisitos con árbol jerárquico.
    \item Apertura, respuesta inicial y estados operativos de No Conformidades.
    \item Marcado de requisitos como No Aplica (NA) con exclusión del cálculo de cumplimiento.
    \item Gestión de acciones correctivas con estructura de lista enlazada, tablero Kanban con drag-and-drop e historial campo por campo.
    \item Verificación y cierre de NC.
    \item Chat contextual implementado con WebSockets para respuestas inmediatas y trazabilidad conversacional.
    \item Motor de novedades y notificaciones apoyado en la misma conexión WebSocket.
    \item Carga, vinculación, revisión, visualización e historial de evidencias almacenadas en Google Drive.
    \item Cálculos de cumplimiento con lógica robusta.
    \item Dashboard analítico con gráficos radar de cumplimiento por cláusula, KPIs de resolución y vistas diferenciadas por rol.
    \item Lobbies diferenciados por rol (Administrador, Evaluador, Responsable SGC, Operativo).
    \item Pruebas unitarias, de integración y de partición de equivalencia.
\end{enumerate}

\subsection{Progreso visual del frontend}

Como evidencia complementaria del avance del MVP, se presentan capturas de las vistas más relevantes del frontend. Estas imágenes se incluyen solo como referencia visual del estado de implementación y no buscan documentar exhaustivamente todas las pantallas del sistema. Se omiten las vistas del rol Evaluador porque, respecto de las vistas del Responsable del SGC, el cambio principal corresponde a las funciones disponibles y permisos habilitados, manteniéndose una estructura visual equivalente.

\begin{figure}[H]
    \centering
    \begin{minipage}{0.48\textwidth}
        \centering
        \includegraphics[width=\textwidth]{frontend/login.png}
        \caption{Vista de inicio de sesión del sistema.}
        \label{fig:frontend-login}
    \end{minipage}
    \hfill
    \begin{minipage}{0.48\textwidth}
        \centering
        \includegraphics[width=\textwidth]{frontend/Dashboard.png}
        \caption{Vista de dashboard principal.}
        \label{fig:frontend-dashboard}
    \end{minipage}
\end{figure}

\begin{figure}[H]
    \centering
    \begin{minipage}{0.48\textwidth}
        \centering
        \includegraphics[width=\textwidth]{frontend/Espacios.png}
        \caption{Vista de gestión de espacios de trabajo.}
        \label{fig:frontend-espacios}
    \end{minipage}
    \hfill
    \begin{minipage}{0.48\textwidth}
        \centering
        \includegraphics[width=\textwidth]{frontend/Usuarios.png}
        \caption{Vista de administración de usuarios.}
        \label{fig:frontend-usuarios}
    \end{minipage}
\end{figure}

\begin{figure}[H]
    \centering
    \begin{minipage}{0.48\textwidth}
        \centering
        \includegraphics[width=\textwidth]{frontend/Evidencias.png}
        \caption{Vista de gestión y revisión de evidencias.}
        \label{fig:frontend-evidencias}
    \end{minipage}
    \hfill
    \begin{minipage}{0.48\textwidth}
        \centering
        \includegraphics[width=\textwidth]{frontend/NC.png}
        \caption{Vista de seguimiento de No Conformidades.}
        \label{fig:frontend-nc}
    \end{minipage}
\end{figure}

\begin{figure}[H]
    \centering
    \includegraphics[width=0.72\textwidth]{frontend/Chat.png}
    \caption{Vista de chat contextual y comunicación operativa.}
    \label{fig:frontend-chat}
\end{figure}


\subsection{Módulos pendientes y riesgos}

Los principales pendientes se concentran en cerrar las correcciones de seguridad IDOR documentadas en la deuda técnica (DT-003), realizar pruebas formales de aislamiento multi-tenancy, ejecutar las pruebas UAT con la contraparte y completar el cierre documental del proyecto.

El riesgo técnico más importante es asegurar que la segregación multi-tenancy funcione correctamente en todos los endpoints, vistas y canales WebSocket. Si bien se realizó una auditoría parcial de endpoints (documentada en DT-003), algunos endpoints aún requieren verificación formal. El segundo riesgo corresponde a la deuda técnica acumulada en componentes frontend de alta complejidad, particularmente \texttt{RequirementContent.jsx}, cuyo tamaño y acoplamiento dificultan el mantenimiento futuro y ya están documentados como DT-001. Las 7 deudas técnicas registradas (DT-001 a DT-008) tienen severidad variable: una crítica (DT-003), cuatro altas y dos medias; la mayoría cuenta con plan de corrección definido para antes de la entrega final.



\section{Validación del trabajo}

La validación del MVP se plantea como un proceso incremental, asociado a los requisitos funcionales vigentes del SRS y a los criterios de aceptación de las historias de usuario que permanecen dentro del alcance luego de los ajustes. La validación no busca demostrar que el producto completo ya está terminado, sino comprobar que el núcleo implementado funciona de forma consistente, trazable y segura para los roles principales definidos.

\subsection{Pruebas unitarias e integración}

Las pruebas unitarias e integración existentes evalúan componentes básicos del MVP y deben interpretarse como una validación inicial, no como una cobertura completa de todos los escenarios funcionales, de seguridad o multi-cliente. Las pruebas unitarias se concentran en controladores y servicios aislados, mientras que las pruebas de integración disponibles verifican flujos base de autenticación, usuarios y servicio de evidencias/Drive. Por ello, aunque la suite ejecutada respalda el funcionamiento inicial de módulos relevantes, todavía no reemplaza las pruebas multi-tenancy ni las pruebas UAT planificadas para el cierre.

La evidencia automatizada actual se organiza en dos carpetas principales. En integración se consideran los archivos \texttt{auth.integration.test.js}, \texttt{driveService.integration.test.js} y \texttt{users.integration.test.js}. En unitarias se consideran \texttt{accionesController.test.js}, \texttt{authController.test.js}, \texttt{chatController.test.js}, \texttt{driveService.unit.test.js}, \texttt{evidenceController.test.js}, \texttt{ncController.test.js}, \texttt{updateAction.equivalence.test.js}, \texttt{updateEvidence.test.js}, \texttt{userController.test.js} y \texttt{workspacesController.test.js}. Estos archivos cubren validaciones básicas de cada tema y sirven como evidencia de avance técnico.

De particular relevancia metodológica es el archivo \texttt{updateAction.equivalence.test.js}, que aplica la técnica de partición de equivalencia para validar la actualización de acciones correctivas. Esta técnica, propia del testing de caja negra, permite dividir el dominio de entrada en clases de equivalencia y seleccionar casos representativos que cubran comportamientos válidos, inválidos y límites. Su inclusión en la suite aporta rigor a la validación y constituye una evidencia concreta de la aplicación de técnicas formales de testing en el proyecto.

Como evidencia de avance, se ejecutó la suite automatizada disponible en el proyecto y se obtuvieron resultados exitosos en las pruebas unitarias y de integración existentes. Las Figuras \ref{fig:pruebas-unitarias-1} y \ref{fig:pruebas-unitarias-2} muestran la salida de terminal correspondiente a una misma ejecución, separada en dos capturas para visualizar la totalidad de los resultados.

\begin{figure}[H]
    \centering
    \includegraphics[width=0.95\textwidth]{img/Pruebas1.png}
    \caption{Evidencia de ejecución exitosa de pruebas unitarias, primera parte de la salida de terminal.}
    \label{fig:pruebas-unitarias-1}
\end{figure}

\begin{figure}[H]
    \centering
    \includegraphics[width=0.95\textwidth]{img/Pruebas2.png}
    \caption{Evidencia de ejecución exitosa de pruebas unitarias e integración, segunda parte de la salida de terminal.}
    \label{fig:pruebas-unitarias-2}
\end{figure}

\subsection{Validación multi-tenancy}

Dado que la plataforma operará con múltiples clientes, la validación de aislamiento multi-tenancy es prioritaria. Esta prueba debe comprobar que un usuario asociado a un espacio de trabajo no pueda consultar, modificar ni visualizar requisitos, evidencias, No Conformidades, acciones correctivas o historiales pertenecientes a otro cliente. Esta dimensión es crítica porque el valor del sistema depende de mantener confidencialidad y separación estricta de datos entre organizaciones.

Se ha realizado una auditoría parcial de endpoints (documentada en DT-003) que identificó vulnerabilidades IDOR en ciertos controladores, de las cuales varias ya fueron corregidas. La validación formal completa de todos los endpoints, vistas y canales WebSocket se encuentra en curso y constituye una de las tareas prioritarias para el cierre del proyecto. Adicionalmente, el documento \texttt{PlanDePrueba.md} en la raíz del repositorio define 9 fases de validación que incluyen una fase específica para pruebas de seguridad multi-tenancy.

\subsection{Pruebas de aceptación de usuario}

Las pruebas UAT se proyectan sobre tres casos representativos: Administrador, Evaluador y Responsable del SGC. Para cada rol se espera verificar las historias de usuario vigentes y sus criterios de aceptación: administración de espacios y usuarios, navegación por requisitos ISO, apertura y seguimiento de No Conformidades, carga y revisión de evidencias, comentarios en tiempo real, notificaciones WebSocket e historial acotado. Estas pruebas aún están pendientes y se realizarán de acuerdo con el cronograma estipulado para el cierre del proyecto. Las historias relacionadas con correo transaccional, S3, carga CSV, edición documental, flujos complejos o soporte multi-norma quedan excluidas de esta validación porque fueron removidas del alcance del MVP.

\begin{table}[H]
\centering
\begin{tabular}{|p{3.2cm}|p{7.2cm}|p{4cm}|}
\hline
\rowcolor{gray!20}\textbf{Tipo de prueba} & \textbf{Propósito} & \textbf{Estado esperado} \\ \hline
Unitarias & Verificar controladores, servicios y funciones aisladas en escenarios básicos. & Ejecutadas exitosamente en los archivos unitarios disponibles; evidencia incorporada en capturas de terminal. \\ \hline
Integración & Comprobar flujos base de autenticación, usuarios y servicio de evidencias/Drive. & Ejecutadas exitosamente en los archivos de integración disponibles; evidencia incorporada en capturas de terminal. \\ \hline
Multi-tenancy & Confirmar aislamiento estricto entre espacios de trabajo. & Pendiente; se ejecutará según cronograma como validación crítica final. \\ \hline
UAT & Validar flujos reales con Administrador, Evaluador y Responsable del SGC. & Pendiente; se preparará, ejecutará y cerrará según cronograma con la contraparte. \\ \hline
\end{tabular}
\caption{Plan de validación del MVP.}
\label{tab:validacion}
\end{table}


\section{Resultados}

Los resultados obtenidos muestran que el núcleo del sistema cuenta con capacidades funcionales completas para ejecutar un flujo GAP integral. La plataforma permite administrar clientes y usuarios, iniciar sesión, acceder a espacios de trabajo, navegar requisitos ISO, marcar requisitos como No Aplica, abrir No Conformidades, registrar respuestas, gestionar acciones correctivas mediante tablero Kanban con estructura padre/hijo, cargar evidencias en Google Drive, visualizar dashboards analíticos con gráficos radar por rol y comunicarse mediante chat y notificaciones en tiempo real vía WebSockets. Esto significa que el proyecto ya superó la etapa de diseño conceptual y posee una base funcional sólida sobre la cual se pueden realizar pruebas de aceptación cercanas al uso real.

El avance cuantitativo debe interpretarse de forma conservadora. Las 245 unidades completadas representan un 87,2\% del total actualizado de 281 unidades, mientras que al considerar las 12 unidades en progreso el avance operativo observado sube a 91,5\%. La diferencia entre ambos indicadores es relevante: el 87,2\% refleja solo trabajo finalizado y verificable, mientras que el 91,5\% muestra el esfuerzo actualmente encaminado hacia cierre. Por esta razón, el indicador conservador es más apropiado para reportar estado real ante la contraparte y el profesor guía.

La distribución de 245 unidades completadas, 12 en progreso y 24 pendientes indica que la mayor parte del núcleo técnico ya fue abordada, quedando pendientes tareas de validación final y cierre documental. Las tareas pendientes asociadas a correcciones IDOR, pruebas multi-tenancy formales, UAT y documentación final no son accesorias: son necesarias para demostrar que el sistema no solo funciona en módulos aislados, sino que también responde a criterios de seguridad, gobernanza y aceptación de usuario.

\begin{table}[H]
\centering
\begin{tabular}{|p{3.4cm}|p{4.1cm}|p{3.2cm}|p{4cm}|}
\hline
\rowcolor{gray!20}\textbf{Validación} & \textbf{Propósito} & \textbf{Estado} & \textbf{Interpretación} \\ \hline
Pruebas unitarias & Verificar controladores, servicios y funciones aisladas en escenarios básicos. & Ejecutadas exitosamente en la suite disponible. & Respaldan estabilidad inicial, pero no implican cobertura total. \\ \hline
Pruebas de integración & Comprobar flujos base de autenticación, usuarios y servicio de evidencias/Drive. & Ejecutadas exitosamente en la suite disponible. & Respaldan integración inicial de módulos específicos, no todos los flujos del sistema. \\ \hline
Multi-tenancy & Verificar separación estricta de datos entre clientes y espacios de trabajo. & Pendiente según cronograma. & Es condición crítica antes de considerar estable la entrega final. \\ \hline
UAT & Validar flujos reales con Administrador, Evaluador y Responsable del SGC. & Pendiente según cronograma. & Permitirá confirmar aceptación funcional con contraparte. \\ \hline
\end{tabular}
\caption{Análisis de resultados de validación del MVP.}
\label{tab:analisis-validacion-resultados}
\end{table}

Desde el punto de vista de ingeniería, la reducción de alcance permitió disminuir dependencias externas y concentrar el esfuerzo en estabilidad funcional. La incorporación de WebSockets para chat en tiempo real y notificaciones internas fue implementada exitosamente, reutilizando una única conexión para ambos propósitos y reduciendo la sobrecarga de mantener canales separados. Esta decisión, que en el informe preliminar figuraba como una reconsideración de la exclusión inicial, hoy constituye un componente estable y funcional del sistema. Se mantienen fuera del alcance el correo transaccional, S3, carga CSV y edición documental.

Comparado con alternativas comerciales como ISWO, GOAudits o GRCTools, el MVP actual todavía es más limitado en amplitud funcional. Sin embargo, la incorporación del dashboard analítico con gráficos radar por cláusula, el tablero Kanban de acciones correctivas y la integración con Google Drive para evidencias reduce significativamente la distancia con plataformas comerciales en cuanto a capacidades de visualización y trazabilidad. La entrevista de validación permite precisar que su valor inicial está en ser simple de utilizar, centralizado, ajustado al flujo real de la consultora y suficientemente comprensible para clientes que no tienen experiencia previa con ISO. Según la evaluación de María, el MVP ya resulta útil porque guarda los datos, organiza los componentes necesarios y permite comenzar a operar sobre el núcleo GAP/NC. La interpretación de los resultados, por tanto, es que el proyecto avanza correctamente como prototipo funcional y base de producto, pero requiere completar pruebas de seguridad, validación formal multi-tenancy y cierre documental antes de considerarse una entrega estable.

Un resultado metodológico relevante es la creación de un registro formal de deuda técnica (8 ítems documentados con severidad, tipo, estado y plan de corrección), así como 11 Architecture Decision Records (ADR) que documentan las decisiones arquitectónicas del proyecto con sus alternativas consideradas y razones de descarte. Ambos artefactos constituyen evidencia de rigor en la gestión del desarrollo y serán analizados en secciones dedicadas más adelante en este informe.

\section{Deuda técnica y lecciones aprendidas}

Durante el desarrollo del MVP se adoptó un enfoque de alta velocidad de implementación que, si bien permitió alcanzar un avance funcional significativo en tiempo reducido, generó deuda técnica en varios frentes. Para gestionar este riesgo de forma transparente, se creó un registro formal de deuda técnica en la carpeta \texttt{docs/deuda-tecnica/}, donde cada ítem está documentado con severidad, tipo, archivos afectados, descripción del problema, impacto y plan de corrección.

\subsection{Registro de deuda técnica}

\begin{table}[H]
\centering
\begin{tabular}{|p{1.8cm}|p{2.2cm}|p{1.5cm}|p{2cm}|p{5.2cm}|}
\hline
\rowcolor{gray!20}\textbf{ID} & \textbf{Título} & \textbf{Severidad} & \textbf{Tipo} & \textbf{Plan de corrección} \\ \hline
DT-001 & \texttt{RequirementContent.jsx} gigante (+1200 líneas) & Media & Mantenibilidad & Refactorizar en componentes más pequeños; dividir lógica de vista de NC, evidencias y evaluaciones. \\ \hline
DT-002 & Nomenclatura NC vs Brecha inconsistente en backend & Baja & Claridad & Se acepta como deuda consciente: el backend mantiene ``NC'' como término canónico, el frontend traduce a ``Brecha''. No se corregirá en esta etapa. \\ \hline
DT-003 & Falta validación IDOR en algunos endpoints & Alta & Seguridad & Auditoría y corrección de endpoints que no validan pertenencia al workspace. Parcialmente avanzado; prioritario para el cierre. \\ \hline
DT-004 & \texttt{console.log} de debug en producción & Baja & Limpieza & Reemplazar por logger con niveles configurable. Corrección simple, planificada para antes de la entrega. \\ \hline
DT-005 & Falta de tests unitarios en controllers críticos & Media & Calidad & Incrementar cobertura de pruebas en controladores de acciones correctivas y evidencias. \\ \hline
DT-007 & Historial de acciones se inserta en dos filas & Baja & Mantenibilidad & Unificar en una sola inserción con snapshot completo. Baja prioridad, no bloquea entrega. \\ \hline
DT-008 & No se valida que acciones correctivas no queden vacías & Media & Validación & Agregar validación en backend y frontend para evitar creación de acciones sin descripción ni responsable. \\ \hline
\end{tabular}
\caption{Registro de deuda técnica del proyecto.}
\label{tab:deuda-tecnica}
\end{table}

\subsection{Lecciones aprendidas sobre el uso de IA generativa}

El proyecto constituye un caso de estudio sobre el uso de IA generativa (LLMs) como herramienta de productividad en el desarrollo de software. La experiencia permitió identificar una distinción relevante entre dos enfoques:

\begin{itemize}
    \item \textbf{``Vibe coding'':} prioriza la velocidad de implementación sobre la revisión. Las funcionalidades se entregan sin revisión inmediata del código generado, las correcciones se postergan para ``un momento futuro'', y los archivos crecen sin refactorización. Este enfoque, utilizado en las primeras etapas del proyecto, generó la mayor parte de la deuda técnica registrada.
    \item \textbf{``Vibe engineering'':} incorpora diseño previo, segmentación granular de tareas, revisión inmediata de cada entrega y tareas de auditoría intercaladas. La regla práctica adoptada es: si una tarea produce más de 50 líneas de código, debió dividirse en subtareas más pequeñas.
\end{itemize}

La transición del primer enfoque al segundo durante el desarrollo del proyecto permitió estabilizar la calidad del código en las etapas finales. Las lecciones extraídas son:

\begin{enumerate}
    \item \textbf{Diseñar y segmentar tareas de forma granular:} cada tarea debe ser lo suficientemente pequeña para que su revisión sea rápida y efectiva (idealmente bajo 50 líneas de código nuevo o modificado).
    \item \textbf{Revisar inmediatamente cada tarea completada:} la revisión inmediata evita la acumulación de defectos que luego requieren sesiones largas de depuración.
    \item \textbf{Intercalar tareas de auditoría/mantenibilidad:} cada 3 a 5 tareas funcionales, dedicar una tarea a revisar calidad, eliminar código muerto, verificar errores de consola y asegurar que no se introdujeron regresiones.
\end{enumerate}

Estas lecciones son particularmente relevantes en el contexto de un taller de titulación, donde el uso de herramientas de IA generativa es cada vez más común y donde la capacidad de integrarlas con rigor metodológico constituye una competencia profesional diferenciadora.


\section{Decisiones de Arquitectura (ADR)}

Como parte del aseguramiento de la calidad del diseño, se documentaron 11 Architecture Decision Records (ADR) en el archivo \texttt{docs/ADR.md}. Cada ADR registra una decisión arquitectónica relevante, las alternativas consideradas, los criterios de evaluación y la justificación de la opción seleccionada. Esta práctica, inspirada en la metodología de Michael Nygard, permite mantener trazabilidad sobre las decisiones que dieron forma al sistema y facilita la incorporación de nuevos desarrolladores o la revisión futura del diseño.

\begin{table}[H]
\centering
\begin{tabular}{|p{0.8cm}|p{4.2cm}|p{5.5cm}|p{4cm}|}
\hline
\rowcolor{gray!20}\textbf{ADR} & \textbf{Decisión} & \textbf{Alternativas consideradas} & \textbf{Criterio principal} \\ \hline
ADR-001 & JWT + refresh token en HTTP-only cookie & Sesiones server-side, OAuth2, Basic Auth & Seguridad stateless compatible con SPA \\ \hline
ADR-002 & MySQL con InnoDB & PostgreSQL, MongoDB, SQLite & Compatibilidad con infraestructura existente del cliente \\ \hline
ADR-003 & Express.js sobre Node.js & Fastify, NestJS, Spring Boot & Curva de aprendizaje y productividad del equipo \\ \hline
ADR-004 & REST API con JSON & GraphQL, gRPC, SOAP & Simplicidad, tooling y ecosistema frontend \\ \hline
ADR-005 & Google Drive como repositorio de evidencias & AWS S3, Azure Blob, MySQL (BLOB) & Costo, simplicidad de integración y segregación por workspace \\ \hline
ADR-006 & WebSockets con librería \texttt{ws} & Socket.io, SSE, polling & Misma conexión para chat + notificaciones, sin dependencias adicionales \\ \hline
ADR-007 & Vite + React SPA & Next.js, Create React App, Angular & Velocidad de desarrollo y ecosistema de componentes \\ \hline
ADR-008 & Multi-tenancy por workspace con middleware de autorización & Esquemas separados por BD, tenant ID en JWT, microservicios por cliente & Simplicidad operacional para MVP \\ \hline
ADR-009 & RBAC con middleware encadenado & ABAC, ReBAC, Oso & Suficiente para 4 roles; extensible a futuro \\ \hline
ADR-010 & Seeding SQL idempotente para estructura ISO & Migraciones programáticas, carga dinámica CSV, API externa & Reproducibilidad y control de versiones \\ \hline
ADR-011 & Estructura de acciones correctivas como linked-list (\texttt{accion\_previa\_id}) & Árbol jerárquico completo, tabla de relaciones muchos-a-muchos & Simple de implementar, suficiente para MVP, extensible \\ \hline
\end{tabular}
\caption{Resumen de las 11 Architecture Decision Records del proyecto.}
\label{tab:adr}
\end{table}

Cada ADR incluye una tabla comparativa detallada con las razones de descarte de cada alternativa no seleccionada, lo que demuestra que las decisiones no fueron arbitrarias sino el resultado de un análisis estructurado de trade-offs. La existencia de este registro constituye un activo valioso para la mantenibilidad futura del sistema y para la comunicación del diseño ante la contraparte técnica.

\section{Plan de cierre}

Para completar el proyecto se propone priorizar las tareas restantes en el siguiente orden:

\begin{enumerate}
    \item Cerrar las correcciones de seguridad IDOR pendientes en los endpoints de chat y acciones correctivas (DT-003), asegurando que todos los controladores verifiquen pertenencia al workspace del usuario autenticado.
    \item Ejecutar pruebas formales de aislamiento multi-tenancy sobre todos los endpoints, vistas y canales WebSocket, verificando que ningún usuario pueda acceder a datos de otro workspace.
    \item Preparar y ejecutar casos UAT representativos para los cuatro roles: Administrador, Evaluador, Responsable del SGC y Operativo.
    \item Refactorizar \texttt{RequirementContent.jsx} (DT-001) para reducir su tamaño y acoplamiento, mejorando la mantenibilidad del frontend.
    \item Completar la cobertura de pruebas unitarias en controladores críticos (DT-005).
    \item Corregir deudas técnicas de severidad baja (DT-004, DT-007, DT-008) y realizar limpieza general de código.
    \item Actualizar capturas de frontend y diagramas C4 con el estado final del sistema.
    \item Cerrar el informe final y preparar la presentación de defensa.
\end{enumerate}

En cuanto al cierre de deuda técnica, se establece el siguiente criterio: los ítems DT-003 (IDOR) y DT-001 (refactoring de RequirementContent) deben ser corregidos antes de la entrega final por su impacto en seguridad y mantenibilidad respectivamente. Los ítems DT-002 (nomenclatura NC vs Brecha) y DT-007 (historial duplicado) se aceptan como deuda consciente que no bloquea la entrega del MVP. Los ítems DT-004, DT-005 y DT-008 se corregirán en la medida que el tiempo lo permita, priorizando DT-005 por su impacto en la confiabilidad de las pruebas.

\newpage

\section{Conclusión}

El proyecto aborda una problemática concreta de Asesorías ISO RCR SpA: la necesidad de digitalizar y estandarizar el diagnóstico GAP, la gestión de evidencias y el seguimiento de No Conformidades para ISO 9001:2015. La solución propuesta tiene una relación directa con la operación de la empresa y con su posibilidad de ofrecer un servicio web escalable a sus clientes.

Respecto del logro de objetivos, el primer objetivo específico se encuentra cubierto mediante el SRS, la definición de roles, el modelo multi-tenancy y la documentación C4 complementada con ADRs. El segundo objetivo muestra avance completo mediante autenticación JWT con refresh token, autorización RBAC con middleware encadenado, CRUD de usuarios y espacios de trabajo, activación de cuenta y timeout por inactividad. El tercer objetivo se aborda mediante el seeding idempotente de requisitos ISO 9001:2015. El cuarto objetivo presenta avance significativo: el motor GAP está funcional con evaluación de cumplimiento, marcado NA, apertura de No Conformidades con máquina de estados completa, acciones correctivas con estructura padre/hijo, tablero Kanban con drag-and-drop y dashboard analítico con gráficos radar por cláusula y KPIs segmentados por rol. El quinto objetivo se cumple mediante la integración con Google Drive para almacenamiento de evidencias con carpetas por workspace, flujo de validación (Pendiente/Aceptado/Rechazado) e historial visual. El sexto objetivo muestra avance parcial: las pruebas unitarias, de integración y de partición de equivalencia fueron ejecutadas exitosamente, las vulnerabilidades IDOR en controladores de chat y acciones fueron corregidas, pero las pruebas formales multi-tenancy y UAT con la contraparte permanecen como tareas de cierre.

Las principales limitaciones del trabajo se relacionan con el alcance del MVP y con el tiempo disponible. La versión actual no incluye recuperación de contraseña por correo, onboarding automatizado, carga dinámica CSV, edición documental en línea, flujos complejos de aprobación, soporte multi-norma ni automatizaciones avanzadas. En contrapartida, el sistema incorpora capacidades no previstas originalmente que agregan valor significativo: WebSockets para chat y notificaciones en tiempo real, dashboard analítico con radares, Kanban de acciones correctivas, almacenamiento en Google Drive y registro formal de deuda técnica y decisiones arquitectónicas. El MVP no alcanza la amplitud funcional de plataformas comerciales existentes; su fortaleza está en ser una base propia, simple de usar, adaptable y rigurosamente alineada al flujo de negocio de la consultora.

Como oportunidades de mejora para la empresa, se identifica la posibilidad de evolucionar el sistema hacia una plataforma diferenciada por cuatro atributos principales: simplicidad de uso, centralización de la información, auditabilidad detallada e inteligencia aplicada. Estas mejoras permitirían reducir trabajo duplicado, evitar el uso de múltiples sistemas, aumentar trazabilidad, proteger mejor la información de clientes, generar nuevos servicios digitales y fortalecer la propuesta comercial de Asesorías ISO RCR SpA.

El MVP es funcional como base operativa completa para el flujo GAP, pero no debe considerarse una entrega final estable hasta que las pruebas multi-tenancy y UAT estén cerradas, las correcciones IDOR estén completamente verificadas y la separación entre datos de clientes sea exhaustivamente validada.

En términos del proceso de desarrollo, el proyecto deja como aprendizaje la importancia del diseño previo y la revisión inmediata del código generado con herramientas de IA. La transición desde un enfoque de ``vibe coding'' hacia ``vibe engineering'' —con tareas granulares, revisión inmediata y auditorías intercaladas— permitió estabilizar la calidad en las etapas finales. La documentación de deuda técnica y ADRs constituye un activo que trasciende el alcance del MVP y sienta las bases para una evolución ordenada del producto.

La proyección post-MVP consiste en consolidar las pruebas de aceptación, cerrar la deuda técnica prioritaria, estabilizar el despliegue y luego priorizar funcionalidades de mayor valor estratégico. En particular, la incorporación progresiva de IA agente, importación inteligente, integración entre componentes, historial detallado y soporte multi-norma permitiría transformar el MVP en una plataforma SaaS más robusta, competitiva y alineada con la visión de largo plazo del proyecto. El plan de pruebas documentado en \texttt{PlanDePrueba.md} proporciona una base metodológica para la validación futura del sistema en cada etapa de evolución.

\section{Referencias}

\subsection{Documentación y evidencias internas}

\begin{enumerate}
    \item Documento de Especificación de Requisitos de Software (SRS), ``Plataforma Web para la Ejecución de GAP Análisis y Gestión de Evidencias por Cláusulas en ISO 9001:2015'', versión 3.1, 10 de abril de 2026.
    \item Hoja de seguimiento ``Avance Proyecto ISO.xlsx'', estado de tareas y horas del proyecto.
    \item Reporte preliminar ``Reporte de Avances Taller de Titulación'', Software web de gestión ISO, 05 de junio de 2026.
    \item Registro de Architecture Decision Records (ADR), 11 decisiones documentadas con alternativas y criterios de selección, 2026.
    \item Registro de Deuda Técnica, 8 ítems documentados con severidad, tipo y plan de corrección, 2026.
    \item Plan de Pruebas E2E, ``PlanDePrueba.md'', 9 fases de validación para el MVP, 2026.
    \item Documento de componentes críticos C4, ``C4-componentes-criticos.md'', diagramas detallados de arquitectura, 2026.
    \item Entrevista de validación estratégica y técnica del SGC con María, consultora especialista en sistemas de gestión ISO, 05 de junio de 2026. Disponible en: \url{https://docs.google.com/document/d/1ei6-ehbhsUr21Wq_-4oHyQDtDc5OqYl41NJ8X2qaqeA/edit?usp=sharing}.
    \item Correos electrónicos de avance semanal del proyecto, ``Informe semanal de progreso -- Proyecto ISO'', comunicaciones de seguimiento, 05 de junio de 2026.
    \item Carpeta de Google Drive con registro de avances del proyecto para revisión del profesor guía. Disponible en: \url{https://drive.google.com/drive/folders/1wh1V29T4QudkEej6gY3Nv1SmFBq_m3dH}.
    \item Correo electrónico con validación de María, consultora especialista en sistemas de gestión ISO, comunicación personal incluida como evidencia de avance, 05 de junio de 2026.
    \item Asesorías ISO RCR SpA. Certificado de aprobación del curso ``Interpretación de ISO 9001:2015'', emitido como apoyo formativo para el desarrollo del proyecto, 2026.
\end{enumerate}

\subsection{Fuentes externas}

\begin{enumerate}
    \item International Organization for Standardization, ISO 9001 Quality Management Systems. Disponible en: \url{https://www.iso.org/iso-9001-quality-management.html}.
    \item React Documentation, componentes y arquitectura de interfaces. Disponible en: \url{https://react.dev/}.
    \item Vite Documentation, herramienta de desarrollo frontend para aplicaciones web. Disponible en: \url{https://vite.dev/}.
    \item Express.js Documentation, framework web para Node.js. Disponible en: \url{https://expressjs.com/}.
    \item Stack Overflow Developer Survey, tendencias de tecnologías web. Disponible en: \url{https://survey.stackoverflow.co/2024/technology/}.
    \item IBM, REST APIs y arquitectura de servicios. Disponible en: \url{https://www.ibm.com/topics/rest-apis}.
    \item Amazon Web Services. Amazon S3, almacenamiento de objetos en la nube. Disponible en: \url{https://aws.amazon.com/es/s3/}.
    \item MySQL Documentation, InnoDB y propiedades ACID. Disponible en: \url{https://dev.mysql.com/doc/refman/8.4/en/mysql-acid.html}.
    \item IETF RFC 7519, JSON Web Token (JWT). Disponible en: \url{https://www.rfc-editor.org/rfc/rfc7519}.
    \item IETF RFC 6455, The WebSocket Protocol. Disponible en: \url{https://www.rfc-editor.org/rfc/rfc6455}.
    \item Atlassian, Kanban y gestión visual del trabajo. Disponible en: \url{https://www.atlassian.com/agile/kanban}.
    \item Brown, Simon. The C4 Model for Visualising Software Architecture. Sitio web oficial. Disponible en: \url{https://c4model.com/}.
    \item ISWO, SIG ISWO Software para sistemas de gestión ISO y HSEQ. Disponible en: \url{https://iswo.com.co/software/}.
    \item GOAudits, software de auditorías e inspecciones. Disponible en: \url{https://goaudits.com/}.
    \item GRCTools, herramientas para gestión GRC. Disponible en: \url{https://grctools.software/}.
\end{enumerate}


\end{document}
